\documentclass[10pt,a4paper]{amsart}
\usepackage[margin=19mm]{geometry}
\usepackage{amsmath,amssymb,mathtools}
\usepackage[scr=rsfs,cal=euler]{mathalfa}
\usepackage{xfrac}
\usepackage[unicode,hidelinks,bookmarks=false]{hyperref}
\newcommand{\ie}{i.e.\ }
\newcommand{\cf}{cf.\ }
\newcommand{\resp}{respectively\ }
\newcommand{\Subsection}{\subsection}
\providecommand{\nobreakdash}{\nobreak\hbox{-}}
\setlength{\emergencystretch}{2em}
\multlinegap0pt
\usepackage[shortlabels]{enumitem}
\usepackage{amssymb}
\usepackage{mathtools}
\newtheorem{theorem}{Theorem}
\newtheorem{lemma}{Lemma}
\newtheorem{proposition}{Proposition}
\newcommand{\Zp}{\mathbb{Z}_p}
\newcommand{\bigo}{O}
\newcommand{\gr}{\operatorname{gr}}
\DeclareMathOperator{\ordop}{ord}
\newcommand{\vp}{\ordop_p}
\title{Coefficient-Level B\"ottcher Theory for Wild Superattracting Germs of Degree \NoCaseChange{\texorpdfstring{$p^e$}{pe}}}
\author{Rufei Ren}
\date{Source: arXiv:2604.02014v3 (CC BY 4.0)}
\begin{document}
\maketitle

\noindent\emph{Source and licence.} Coefficient-Level B\"ottcher Theory for Wild Superattracting Germs of Degree \NoCaseChange{\texorpdfstring{$p^e$}{pe}}. Version arXiv:2604.02014v3 (CC BY 4.0), distributed by arXiv under the Creative Commons Attribution 4.0 International licence, http://creativecommons.org/licenses/by/4.0/, as recorded in the arXiv licence metadata for this exact version and shown on the abstract page https://arxiv.org/abs/2604.02014v3. Full licence notice: \texttt{licenses/arxiv-2604.02014-CC-BY-4.0.txt}.\par\smallskip

\noindent\textit{Editorial scope.} Counted unit: the article's theorem thm:higher-structure (source0.tex lines 195--231). The article's complete original proof of it is delivered with it (source0.tex lines 1583--1778, one \texttt{proof} environment of 196 lines, which the article places away from the statement). No mathematics is written, repaired or re-derived: the statement and the proof are the article's own lines, in the article's own notation, and no retained line is substituted or re-flowed.  Retained uncounted as the source-local context the proof needs: lines 449--455; lines 476--482; lines 497--506; lines 1084--1086; lines 1088--1090; lines 1138--1150; lines 1235--1258; lines 1299--1312; lines 1475--1502; lines 1520--1550.  Nothing was omitted from any retained span, so the package carries no omission record.  No retained line contains a citation, so the wrapper carries no bibliography and no bibliography entry.  No retained line contains a figure environment, an image-inclusion command or drawing code, so no image is delivered and the package declares no asset dependency.  Editorial formatting only, in addition: an AMS \texttt{amsart} preamble that restates the article's own declarations and macros used by the retained lines, namely the environments \texttt{theorem}, \texttt{lemma}, \texttt{proposition} on separate counters, together with the standard packages those lines need.  The retained environments print in this document as Theorem 1, Lemma 1, Proposition 1 in order of appearance, whereas the article prints the same items with the numbering of its own full document.  The complete native source of the article as distributed in the arXiv e-print for this exact version is bundled unchanged as \texttt{sources/197700/source0.tex}.
\begin{lemma}[Global $B$-coefficient integrality]\label{lem:B-integrality}
For every $k\ge1$ and every monomial coefficient $\gamma_B(\eta)$ occurring in $B_k[x^k]$,
\[
        c(\eta)+d(\eta)\ge e.
\]
Equivalently, $\gamma_B(\eta)\in\Zp$.
\end{lemma}

\begin{proposition}[Unit $B$-terms on divisible classes]\label{prop:B-unit-divisible}
Assume $p\mid k$ and let a monomial in $B_k[x^k]$ have $p$-adic unit scalar coefficient.  If some occurring positive index is not divisible by $p$, then necessarily $k=p$ and the monomial is
\[
        a_0^{q-p}a_1^p.
\]
Consequently, if $p\mid k$ and $k>p$, every unit-coefficient monomial in $B_k[x^k]$ uses only positive indices divisible by $p$.
\end{proposition}

\begin{proposition}[A global $C$-term lemma]\label{prop:C-global}
For every $k\ge1$, every monomial coefficient of $C_k[x^k]$ is $p$-integral.  Moreover:
\begin{enumerate}[label=\textup{(\roman*)},leftmargin=2.2em]
\item if $p\mid k$, then every monomial coefficient is divisible by $p$, so $C_k[x^k]\equiv0\pmod p$;
\item if $k\equiv a\pmod p$ with $1\le a\le p-1$, then the unique surviving term modulo $p$ is $a\,a_{k-1}$, i.e.
\[
        C_k[x^k]\equiv a\,a_{k-1}\pmod p.
\]
\end{enumerate}
\end{proposition}

\begin{equation}\label{eq:a1-r}
        a_1(r)=-p^r,
\end{equation}

\begin{equation}\label{eq:ABC-r}
        a_k(r)=A_k[x^k]-B_k[x^k]-p^rC_k[x^k],
\end{equation}

\begin{lemma}[The first two pure powers]\label{lem:first-two-pure-pe}
For $r\ge1$ and $e\ge2$,
\[
        v_0(r,e)=r,
        \qquad
        v_1(r,e)=pr,
\]
and
\[
        p^{-v_i(r,e)}a_{p^i}(r,e)\equiv -1\pmod p
        \qquad (i=0,1).
\]
\end{lemma}

\begin{proposition}[Exact decomposition on pure powers]\label{prop:pure-decomp}
For every $n\ge1$ there exist polynomials $R_n,S_n\in\Zp[a_1(r),\ldots,a_{p^n-1}(r)]$ such that
\begin{equation}\label{eq:pure-decomp-small}
        a_{p^n}(r)=-\gamma_{n,e}a_{p^{n-1}}(r)^p+pR_n+p^{r+1}S_n
        \qquad(1\le n<e),
\end{equation}
and
\begin{equation}\label{eq:pure-decomp-large}
        a_{p^n}(r)=\alpha_{n,e}a_{p^{n-e}}(r)-\gamma_{n,e}a_{p^{n-1}}(r)^p+pR_n+p^{r+1}S_n
        \qquad(n\ge e).
\end{equation}
Moreover,
\begin{itemize}[leftmargin=2em]
\item every monomial of $R_n$ is a product of lower coefficients whose weighted total index is $p^n$;
\item every monomial of $S_n$ is a product of lower coefficients whose weighted total index is $p^n-1$;
\item all scalar coefficients of $R_n$ and $S_n$ are $p$-integral;
\item
\[
        \vp(\alpha_{n,e})=\Delta_{n,e}=\mu_ep^{n-e}-e,
        \qquad
        p^{-\Delta_{n,e}}\alpha_{n,e}\equiv(-1)^e\pmod p.
\]
\end{itemize}
\end{proposition}

\begin{lemma}[Top-digit lower bound from lower levels]\label{lem:top-digit-lower-pe}
Assume that
\[
        v_{i+1}\le p v_i\qquad (0\le i\le n-2).
\]
Let $e_j$ be nonnegative integers with $0\le j<p^n$ and
\[
        \sum_{j=0}^{p^n-1}e_jj=p^n.
\]
Then
\[
        \sum_{j=0}^{p^n-1}e_j\Lambda(j)\ge p\,v_{n-1}.
\]
\end{lemma}

\begin{proposition}[$\pi$-free initial form of the divisible $B$-term]\label{prop:initial-form-B-pe}
Fix $r\ge1$ and $e\ge2$.  Let $M\ge1$ be not a power of $p$, write
\[
        M=M_0+M_1p+\cdots+M_sp^s,
        \qquad
        0\le M_i\le p-1,
\]
and set $k=pM$.  Assume that
\[
        v_{j+2}(r,e)\le pv_{j+1}(r,e)
        \qquad (0\le j\le s-1).
\]
Define
\[
        Y^{[i]}:=\prod_{j\ge0}Y_j^{i_j}
        \qquad\text{for }i=\sum_{j\ge0} i_jp^j,
\]
so in particular $Y^{[0]}=1$, and set
\[
        \mathcal B_{M,r,e}(Y):=\frac{(pM)!}{p^e}[y^M]
        \left(\sum_{i=0}^{M-1}\frac{Y^{[i]}}{(pi)!}y^i\right)^{p^e}.
\]
Then the degree-$\Lambda(k)$ class of $\mathcal B_{M,r,e}(Y)$ in
$(\gr_\Lambda\mathcal R_{r,e})/(\pi)$ is
\[
        -Y^{[M]}.
\]
\end{proposition}

\begin{lemma}[Abstract pure-power branch pattern and no tie]
\label{lem:abstract-branch-pe}
Fix integers $r\ge1$ and $e\ge2$.  Let $(w_n)_{n\ge0}$ be defined by
\[
        w_0=r,
        \qquad
        w_n=pw_{n-1}\quad(1\le n<e),
\]
and, for $n\ge e$,
\[
        w_n=\min\{\Delta_{n,e}+w_{n-e},\,p w_{n-1}\}.
\]
Set
\[
        s:=\left\lceil\frac r e\right\rceil.
\]
Then the branch word for this abstract recursion is
\[
        (B^{e-1}A)^sB^\infty.
\]
Equivalently, the $A$-branch occurs exactly at $n=e,2e,\ldots,se$, every other level is $B$-dominated, and for every $n\ge e$ one has
\[
        \Delta_{n,e}+w_{n-e}\ne p w_{n-1}.
\]
Moreover
\[
        w_{je+t}=p^t\left(r+\frac{p^{j e}-1}{p-1}-e\,j\right)
        \qquad(0\le j\le s,\ 0\le t\le e-1),
\]
and $w_n=p^{n-e s}w_{e s}$ for all $n\ge e s$.
\end{lemma}

\begin{theorem}[Recursive structure in degree $p^e$]\label{thm:higher-structure}
For every fixed $r\ge 1$ and $e\ge 2$, the following assertions hold.
\begin{enumerate}[label=\textup{(\arabic*)},leftmargin=2.2em]
\item \emph{Global digit-weight lower bound.}  For every $k\ge 1$,
\[
        \vp\bigl(a_k(r,e)\bigr)\ge \Lambda_{r,e}(k).
\]
\item \emph{Leading monomial on divisible non-pure classes.}  If $p\mid k$ and $k$ is not a power of $p$, then
\[
        a_k(r,e)\equiv m_{r,e}(k)\pmod {p^{\Lambda_{r,e}(k)+1}}.
\]
\item \emph{Pure-power recursion.}  One has $v_0(r,e)=r$ and, for $1\le n<e$,
\[
        v_n(r,e)=p\,v_{n-1}(r,e)=p^nr.
\]
For every $n\ge e$,
\[
        v_n(r,e)=\min\{A_{n,e}(r),B_{n,e}(r)\},
        \qquad
        A_{n,e}(r)\ne B_{n,e}(r).
\]
Moreover the corresponding leading branch is
\[
        a_{p^n}(r,e)=
        \begin{cases}
        -\gamma_{n,e}a_{p^{n-1}}(r,e)^p+\bigo\bigl(p^{v_n(r,e)+1}\bigr),&1\le n<e,\\[1mm]
        \alpha_{n,e}a_{p^{n-e}}(r,e)+\bigo\bigl(p^{v_n(r,e)+1}\bigr),& n\ge e,\ A_{n,e}(r)<B_{n,e}(r),\\[1mm]
        -\gamma_{n,e}a_{p^{n-1}}(r,e)^p+\bigo\bigl(p^{v_n(r,e)+1}\bigr),& n\ge e,\ B_{n,e}(r)<A_{n,e}(r).
        \end{cases}
\]
In particular $v_n(r,e)\le p\,v_{n-1}(r,e)$ for all $n\ge 1$.
\item \emph{Subadditivity of the digit weight.}  For every finite sum $N=n_1+\cdots+n_t$ of nonnegative integers,
\[
        \Lambda_{r,e}(N)\le \Lambda_{r,e}(n_1)+\cdots+\Lambda_{r,e}(n_t).
\]
\end{enumerate}
\end{theorem}

\begin{proof}[Proof of Theorem~\ref{thm:higher-structure}]
We verify the induction package $\mathcal I(N)$ by strong induction on $N$.

The base case $\mathcal I(1)$ follows from \eqref{eq:a1-r}: one has
\[
        a_1(r)=-p^r,
        \qquad
        \Lambda(1)=v_0=r,
\]
so (L) holds at $k=1$.  The clauses (D) and (P) are vacuous for $N=1$, and (S) is tautological for totals $M\le1$.

\smallskip
\noindent\emph{Pure powers.}
Let $(w_n)_{n\ge0}$ be the abstract sequence from Lemma~\ref{lem:abstract-branch-pe}.  As part of the induction on $\mathcal I(N)$, whenever a pure-power level $p^j$ has already been treated we also record the equality $v_j=w_j$.

We treat $n=1$ first.  By Lemma~\ref{lem:first-two-pure-pe} we have $v_0=w_0=r$ and $v_1=w_1=pr$.  Now fix $n\ge2$.  Assume (L) and (D) hold for every $k<p^n$, and (P) holds for every pure-power level below $p^n$.  Then (S) is available for all integers whose highest base-$p$ digit is at most $p^{n-1}$: all pure-power levels below $p^n$ have already been treated, so the inequalities $v_{j+1}\le p v_j$ are known for $0\le j\le n-2$, and the carry argument gives subadditivity for every sum supported in those digit places.

We begin with the term $pR_n$ in Proposition~\ref{prop:pure-decomp}.  Every monomial in $R_n$ is a product of lower coefficients whose total index is $p^n$.  By the lower-bound induction hypothesis and Lemma~\ref{lem:top-digit-lower-pe}, its product valuation is at least $p v_{n-1}$.  The extra factor $p$ in $pR_n$ raises the valuation to at least $p v_{n-1}+1$.

For $p^{r+1}S_n$, every monomial in $S_n$ is a product of lower coefficients whose total index is $p^n-1$.  Hence (L) for lower indices and (S) at lower levels show that every monomial contribution to $p^{r+1}S_n$ has valuation at least
\[
        r+1+\Lambda(p^n-1).
\]
Now
\[
        \Lambda(p^n-1)=(p-1)\sum_{i=0}^{n-1}v_i,
\]
because the base-$p$ expansion of $p^n-1$ has all digits equal to $p-1$ below level $n$.  Using $v_{j+1}\le p v_j$ for $0\le j\le n-2$, we obtain
\[
        \sum_{i=0}^{n-2}(p-1)v_i\ge \sum_{i=0}^{n-2}(v_{i+1}-v_i)=v_{n-1}-v_0.
\]
Therefore
\[
\begin{aligned}
        r+1+\Lambda(p^n-1)
        &=1+v_0+(p-1)\sum_{i=0}^{n-1}v_i\\
        &\ge 1+v_0+v_{n-1}-v_0+(p-1)v_{n-1}\\
        &=1+p v_{n-1}.
\end{aligned}
\]
Hence every contribution in $p^{r+1}S_n$ has valuation strictly larger than $p v_{n-1}$.

If $1<n<e$, Proposition~\ref{prop:pure-decomp} gives
\[
        a_{p^n}(r)=-\gamma_{n,e}a_{p^{n-1}}(r)^p+pR_n+p^{r+1}S_n.
\]
The two remainder estimates above show that
\[
        a_{p^n}(r)=-\gamma_{n,e}a_{p^{n-1}}(r)^p+\bigo\bigl(p^{p v_{n-1}+1}\bigr),
\]
so
\[
        v_n=p v_{n-1}=p w_{n-1}=w_n.
\]
This gives the case $n<e$ in (P).

Now assume $n\ge e$.  Substituting the same remainder bounds into Proposition~\ref{prop:pure-decomp}, we obtain
\[
        a_{p^n}(r)=\alpha_{n,e}a_{p^{n-e}}(r)-\gamma_{n,e}a_{p^{n-1}}(r)^p+\bigo\bigl(p^{p v_{n-1}+1}\bigr).
\]
By the already-treated pure-power levels, we have $v_{n-e}=w_{n-e}$ and $v_{n-1}=w_{n-1}$.  Therefore the two main terms have valuations
\[
        A_{n,e}=\Delta_{n,e}+v_{n-e}=\Delta_{n,e}+w_{n-e},
        \qquad
        B_{n,e}=p v_{n-1}=p w_{n-1}.
\]
By Lemma~\ref{lem:abstract-branch-pe}, these two valuations are never equal.  The no-tie assertion is essential here: because the two candidate leading terms have distinct $p$-adic valuations, no cancellation can raise the valuation of $a_{p^n}(r,e)$.  Since both remainder terms have valuation strictly larger than $B_{n,e}$, the dominant term is exactly the smaller of the $A$- and $B$-branches.  Therefore
\[
        v_n=\min\{A_{n,e},B_{n,e}\}=w_n.
\]
If $A_{n,e}<B_{n,e}$, the leading term comes from the $A$-branch and
\[
        a_{p^n}(r)=\alpha_{n,e}a_{p^{n-e}}(r)+\bigo\bigl(p^{v_n+1}\bigr).
\]
If $B_{n,e}<A_{n,e}$, the leading term comes from the $B$-branch and
\[
        a_{p^n}(r)=-\gamma_{n,e}a_{p^{n-1}}(r)^p+\bigo\bigl(p^{v_n+1}\bigr).
\]
Thus the pure-power clause (P) is proved at level $p^n$, and in all cases $v_n\le p v_{n-1}$.

Once these inequalities hold up to level $n$, the carry argument gives (S) for all integers whose base-$p$ expansion uses only the digits $1,p,\ldots,p^n$: replacing $p$ copies of $p^j$ by one copy of $p^{j+1}$ can only decrease or preserve the total weight because $v_{j+1}\le p v_j$.

\smallskip
\noindent\emph{Divisible non-pure indices.}
Let $k=pM$ with $M\ge1$ and $k$ not a power of $p$.  Write
\[
        M=M_0+M_1p+\cdots+M_sp^s,
        \qquad
        0\le M_i\le p-1.
\]
Assume (L) and (D) hold for every $n'<k$, and (P) holds for every pure-power level below $k$.  Then (S) is available for all integers $\le k$: all pure-power levels $p^j\le k$ have already been treated, so the inequalities $v_{j+1}\le p v_j$ are known up to the relevant level, and the carry argument therefore gives subadditivity for every sum with total at most $k$.

We have the following shift inequality.  If $M'\ge1$ and $p^eM'\le k$, then
\begin{equation}\label{eq:shift-ineq-pe}
        \Lambda(p^eM')\le \mu_eM'-e+\Lambda(M'),
\end{equation}
and the inequality is strict if $M'$ is not a power of $p$.  Indeed, writing $M'=\sum_i M'_ip^i$ and using the already-proved pure-power recurrence gives
\[
\begin{aligned}
        \Lambda(p^eM')
        &=\sum_i M_i'v_{i+e}\\
        &\le \sum_i M_i'\bigl(\mu_ep^i-e+v_i\bigr)\\
        &=\mu_eM'-e\sum_i M_i'+\Lambda(M').
\end{aligned}
\]
Since $M'\ge1$, one has $\sum_i M_i'\ge1$, which proves \eqref{eq:shift-ineq-pe}.  If $M'$ is not a power of $p$, then $\sum_i M_i'\ge2$, so the inequality is strict.

If $q\nmid k$, then $A_k[x^k]=0$.  If $k=qM'$, then $M'$ is not a power of $p$, so \eqref{eq:shift-ineq-pe} is strict.  Hence
\[
        \vp\bigl(A_k[x^k]\bigr)\ge \mu_eM'-e+\Lambda(M')>\Lambda(k).
\]
For the $C$-term, every monomial coefficient is divisible by $p$ because $p\mid k$; see Proposition~\ref{prop:C-global}.  Therefore
\[
        \vp\bigl(p^rC_k[x^k]\bigr)\ge 1+r+\Lambda(k-1)>\Lambda(k)
\]
by (S), since $k=(k-1)+1$ and $\Lambda(1)=r$.  Thus
\begin{equation}\label{eq:D-first-reduction-pe}
        a_k(r)\equiv -B_k[x^k]\pmod{p^{\Lambda(k)+1}}.
\end{equation}

If a monomial in $B_k[x^k]$ contains some positive index not divisible by $p$, then, because $k\ne p$, Proposition~\ref{prop:B-unit-divisible} implies that its coefficient is divisible by $p$.  The induction hypothesis (L) and (S) give product valuation at least $\Lambda(k)$, and the coefficient contributes one extra factor of $p$.  Hence modulo $p^{\Lambda(k)+1}$ one may keep only the divisible truncation and write
\begin{equation}\label{eq:divisible-truncation-pe}
        B_k[x^k]\equiv
        \frac{k!}{p^e}[y^M]\left(\sum_{i=0}^{M-1}\frac{a_{pi}(r)}{(pi)!}y^i\right)^{p^e}
        \pmod{p^{\Lambda(k)+1}}.
\end{equation}
Here the index $i=0$ is harmless: both sides are $a_0(r)=m(0)=1$.  For $1\le i<M$, either $pi$ is a pure power, in which case $a_{pi}(r)=m(pi)$ by definition, or $pi$ is divisible and non-pure, in which case the induction hypothesis gives
\[
        a_{pi}(r)\equiv m(pi)\pmod{p^{\Lambda(pi)+1}}.
\]
Consider a monomial in the expansion of \eqref{eq:divisible-truncation-pe} in which at least one factor $a_{pi}(r)$ is replaced by its error term of valuation at least $\Lambda(pi)+1$.  Every remaining factor has valuation at least its digit weight $\Lambda(pj)$, and the scalar coefficient is $p$-integral by Lemma~\ref{lem:B-integrality}.  If the corresponding multiplicities are $\eta_i$, then subadditivity gives
\[
        1+\sum_i \eta_i\Lambda(pi)\ge 1+\Lambda\!\left(\sum_i \eta_i pi\right)=1+\Lambda(k),
\]
so the whole contribution is $\bigo\bigl(p^{\Lambda(k)+1}\bigr)$.  Therefore
\begin{equation}\label{eq:replace-by-mr-pe}
        B_k[x^k]\equiv
        \frac{k!}{p^e}[y^M]\left(\sum_{i=0}^{M-1}\frac{m(pi)}{(pi)!}y^i\right)^{p^e}
        \pmod{p^{\Lambda(k)+1}}.
\end{equation}
The already-proved pure-power part of the induction gives
\[
        v_{j+2}\le pv_{j+1}
        \qquad (0\le j\le s-1),
\]
so Proposition~\ref{prop:initial-form-B-pe} applies to the polynomial $\mathcal B_{M,r,e}(Y)$.  The right-hand side of \eqref{eq:replace-by-mr-pe} is obtained from $\mathcal B_{M,r,e}(Y)$ by the filtered specialization
\[
        Y_j\longmapsto a_{p^{j+1}}(r)
        \qquad (j\ge0),
\]
for which $\vp(a_{p^{j+1}}(r))=v_{j+1}$ by definition.  Hence a term of filtered degree $>\Lambda(k)$ specializes to $\bigo\bigl(p^{\Lambda(k)+1}\bigr)$, while the degree-$\Lambda(k)$ class $-Y^{[M]}$ specializes to $-m(k)$.  Therefore
\[
        B_k[x^k]\equiv -m(k)\pmod{p^{\Lambda(k)+1}}.
\]
Combining this with \eqref{eq:D-first-reduction-pe} gives
\[
        a_k(r)\equiv m(k)\pmod{p^{\Lambda(k)+1}},
\]
which proves (D) at the index $k$.

\smallskip
\noindent\emph{The global lower bound.}
Fix $k\ge1$.  Assume (L) holds for every $n'<k$, (D) holds for every divisible non-pure $n'<k$, and (P) holds for every pure-power level $p^j\le k$.  Then (S) is available for all integers $\le k$: all relevant pure-power levels have already been treated, so the inequalities $v_{j+1}\le p v_j$ are known up to the required height, and the carry argument gives subadditivity for every sum with total at most $k$.

If $k=p^n$ is a pure power, then (P) gives
\[
        \vp\bigl(a_{p^n}(r)\bigr)=v_n=\Lambda(p^n),
\]
so the lower bound holds with equality.  We may therefore assume that $k$ is not a pure power.  We use the recursion \eqref{eq:ABC-r}.

If $q\nmid k$, then $A_k[x^k]=0$.  If $k=qM$, then by the induction hypothesis and \eqref{eq:shift-ineq-pe},
\[
        \vp\bigl(A_k[x^k]\bigr)\ge \mu_eM-e+\Lambda(M)\ge \Lambda(qM)=\Lambda(k).
\]

Take a monomial $\gamma_B(\eta)a^\eta$ in $B_k[x^k]$.  By Lemma~\ref{lem:B-integrality}, $\gamma_B(\eta)\in\Zp$.  By the induction hypothesis,
\[
        \vp(a^\eta)\ge \sum_{n=0}^{k-1}\eta_n\Lambda(n).
\]
Since $\sum_{n=0}^{k-1}\eta_n n=k$, (S) gives
\[
        \sum_{n=0}^{k-1}\eta_n\Lambda(n)\ge \Lambda(k).
\]
Hence every $B$-monomial has valuation at least $\Lambda(k)$.

Take a monomial $\gamma_C(\eta)a^\eta$ in $C_k[x^k]$.  By Proposition~\ref{prop:C-global}, the coefficient $\gamma_C(\eta)$ is $p$-integral, and by the induction hypothesis,
\[
        \vp(a^\eta)\ge \sum_{n=0}^{k-1}\eta_n\Lambda(n)\ge \Lambda(k-1),
\]
because $\sum_{n=0}^{k-1}\eta_n n=k-1$.  Therefore
\[
        \vp\bigl(p^r\gamma_C(\eta)a^\eta\bigr)
        \ge r+\Lambda(k-1)\ge \Lambda(k),
\]
where the last inequality is the special case $k=(k-1)+1$ of (S).  Since every summand in \eqref{eq:ABC-r} has valuation at least $\Lambda(k)$, so does $a_k(r)$.  This gives (L) and finishes the induction proof of Theorem~\ref{thm:higher-structure}.
\end{proof}

\end{document}
